\label{sec:commericial-demonstrations}

This section analyzes the scientific documentation of the
current state of the GFM industry's commercial development,
with the objective of illuminating the degree to which
the control strategies and system models of sections
\ref{sec:math-models} and \ref{sec:control-strategies}
have been tested and validated on physical systems rather
than simply simulation.

\subsection{Distributed IBRs}
The IBRs which have been deployed to-date within
distributed (rooftop) PV \cite{noauthor_grid_2021} and WTG \cite{nguyen_grid-forming_2022}
have been primarily GFL technology. As discussed in section \ref{sec:grid-integration}
converting a GFL IBR to a GFM IBR requires upgrading
the hardware itself, along with a high capital cost. The projected
increase in GFL penetration makes the understanding the relationship
between GFL and GFM dynamics an even more pressing topic.

\subsection{Microgrids}
The Maui power system is a commonly used example of a MG which has 
high instantaneous power of IBRs \cite{kenyon_criticality_2023}.
Other small-scale grids (e.g. grids on small islands) have demonstrated the technical
feasibility of GFM technology \cite{lin_research_2020} while simultaneously presenting 
discrepancies from existing simulation results \cite{kenyon_criticality_2023}.

\subsection{Utility/Grid-scale Implementations}
The authors of \cite{musca_grid-forming_2022}
provide a recent overview of grid-scale systems which have implemented GFM
technology. The reviewed GFM pilot projects typically control a "hybrid power plant,"
consisting of combined solar, wind and battery resources, and are integrated with the grid
on the transmission-level.

The 14 project highlighted by this paper were all completed after 2012, demonstrating
the novelty of large-scale GFM installations. The authors of \cite{musca_grid-forming_2022}
also note that pilot projects represent a valuable source of data for validating
the fidelity of GFM modeling studies.